\documentclass[10pt,a4paper]{amsart}
\usepackage[margin=19mm]{geometry}
\usepackage{amsmath,amssymb,mathtools}
\usepackage[scr=rsfs,cal=euler]{mathalfa}
\usepackage{xfrac}
\usepackage[unicode,hidelinks,bookmarks=false]{hyperref}
\newcommand{\ie}{i.e.\ }
\newcommand{\cf}{cf.\ }
\newcommand{\resp}{respectively\ }
\newcommand{\Subsection}{\subsection}
\providecommand{\nobreakdash}{\nobreak\hbox{-}}
\setlength{\emergencystretch}{2em}
\multlinegap0pt
\usepackage{amssymb}
\usepackage{xcolor}
\newtheorem{claim}{Claim}
\newcommand{\bH}{\mathbb H}
\newcommand{\bR}{\mathbb R}
\title{When a meromorphic function that omits three values is of bounded type}
\author{Alexandre Eremenko, Aleksei Kulikov, Mikhail Sodin}
\date{Source: arXiv:2604.06136v2 (CC BY 4.0)}
\begin{document}
\maketitle

\noindent\emph{Source and licence.} When a meromorphic function that omits three values is of bounded type. Version arXiv:2604.06136v2 (CC BY 4.0), distributed by arXiv under the Creative Commons Attribution 4.0 International licence, http://creativecommons.org/licenses/by/4.0/, as recorded in the arXiv licence metadata for this exact version and shown on the abstract page https://arxiv.org/abs/2604.06136v2. Full licence notice: \texttt{licenses/arxiv-2604.06136-CC-BY-4.0.txt}.\par\smallskip

\noindent\textit{Editorial scope.} Counted unit: the article's claim claim:n (source0.tex lines 811--817). The article's complete original proof of it is delivered with it (source0.tex lines 818--838, one \texttt{proof} environment of 21 lines). No mathematics is written, repaired or re-derived: the statement and the proof are the article's own lines, in the article's own notation, and no retained line is substituted or re-flowed.  No further source-local context is needed: every cross-reference of every retained line resolves inside the two delivered spans.  Nothing was omitted from any retained span, so the package carries no omission record.  No retained line contains a citation, so the wrapper carries no bibliography and no bibliography entry.  No retained line contains a figure environment, an image-inclusion command or drawing code, so no image is delivered and the package declares no asset dependency.  Editorial formatting only, in addition: an AMS \texttt{amsart} preamble that restates the article's own declarations and macros used by the retained lines, namely the environments \texttt{claim} on separate counters, together with the standard packages those lines need.  The retained environments print in this document as Claim 1 in order of appearance, whereas the article prints the same items with the numbering of its own full document.  The complete native source of the article as distributed in the arXiv e-print for this exact version is bundled unchanged as \texttt{sources/198005/source0.tex}.
\begin{claim}\label{claim:n}
For $x\in \bR$, 
\begin{equation}\label{eq:n} 
n(x) \le Cm(x/C)
\end{equation}
with $C= \| W' \|_\infty =\sup\{| W'(\zeta)|\colon \zeta\in\overline\bH(-m)\}$.
\end{claim}

\begin{proof}
Since both functions $m$ and $n$ are even, it suffices to check \eqref{eq:n} 
only for $x\ge 0$. Let, as above, $W=U+iV$. Then $V(t)=n(U(t))$, $t\in\bR$. Since
\[
V(t)={\rm Im\,} W(t) = {\rm Im\,}[ W(t)-W(t-im(t) ],
\]
we have
\[
V(t) \le | W(t) - W(t-im(t))| \le C m(t).
\]
Similarly, 
\[
U(t) = {\rm Re\,} W(t) = {\rm Re\,}[W(t)-W(0)],
\] 
whence, for $t\ge 0$, $U(t)\le |W(t)-W(0)|\le Ct$, 
and therefore, $m(t) \le m(U(t)/C)$. Combining both estimates, we obtain
\[
n(U(t)) = V(t) \le C m(t) \le Cm(U(t)/C).
\] 
Letting $x=U(t)$, we get the claim.
\end{proof}

\end{document}
